\subsection{\texorpdfstring{LOCAL DEFINITION OF A CONTINUOUS HOMOMORPHISM OF \(R^{n}\) INTO A TOPOLOGICAL GROUP}{LOCAL DEFINITION OF A CONTINUOUS HOMOMORPHISM OF R TO THE N INTO A TOPOLOGICAL GROUP}}

Proposition 6 of Chapter V, § 1, no. 4 may be generalized to all the groups \(\mathbf{R}^n\) .

PROPOSITION 2. Let A be a parallelotope in \(\mathbf{R}^n\) which contains o; and let \(f\) be a continuous mapping of A into a topological group G (written multiplicatively) such that \(f(x + y) = f(x)f(y)\) for each pair of points \(x, y\) such that \(x \in A\) , \(y \in A\) and \(x + y \in A\) . Then there exists a unique continuous homomorphism of \(\mathbf{R}^n\) into G which extends \(f\) .

By the same reasoning as in Proposition 6 of Chapter V, §1, no. 4 we show that the homomorphism extending \(f\) , if it exists, is unique. Next, the subgroup \(G_1\) of \(G\) , generated by \(f(A)\) , is commutative; for if \(x\) and y are any two points of A, then \(\frac{1}{2}x\) , \(\frac{1}{2}y\) and \(\frac{1}{2}(x+y)\) belong to A, therefore

\[
f \left(\frac {1}{2} (\boldsymbol {x} + \boldsymbol {y})\right) = f \left(\frac {1}{2} \boldsymbol {x}\right) f \left(\frac {1}{2} \boldsymbol {y}\right) = f \left(\frac {1}{2} \boldsymbol {y}\right) f \left(\frac {1}{2} \boldsymbol {x}\right),
\]

which shows that \(f(\frac{1}{2} x)\) and \(f(\frac{1}{2} y)\) commute; therefore so do \(f(x) = (f(\frac{1}{2} x))^2\) and \(f(y) = (f(\frac{1}{2} y))^2\) , which proves the assertion.

Let \(a_1, a_2, \ldots, a_n\) be \(n\) non-zero vectors contained in \(A\) and proportional to the basis vectors of the parallelotope, and for each index \(i\) let \(D_i\) be the line through \(o\) and \(a_i\) , that is, the set of points \(ta_i\) as \(t\) runs through \(R\) . Let \(A_i\) be the set of all \(t \in R\) such that \(ta_i \in A\) ; then \(A_i\) is an interval of \(R\) containing \([o, i]\) , and the function

\[
f _ {i} (t) = f \left(t \boldsymbol {a} _ {i}\right)
\]

is defined and continuous on \(A_{i}\) and satisfies the relation

\[
f _ {i} (t + t ^ {\prime}) = f _ {i} (t) f _ {i} (t ^ {\prime})
\]

whenever \(t, t'\) and \(t + t'\) all belong to \(A_i\) . By Proposition 6 of Chapter V, §1, no. 4, there exists a continuous homomorphism \(\overline{f}_i\) of \(\mathbf{R}\) into \(G\) which extends \(f_i\) . Since \(\mathbf{R}^n\) is the direct sum of the subgroups \(D_i\) , we can define a homomorphism \(\overline{f}\) of \(\mathbf{R}^n\) into the abelian group \(G_1\) by the rule \(\overline{f}(x) = \prod_{i=1}^{n} \overline{f}_i(t_i)\) where \(x = \sum_{i=1}^{n} t_i a_i\) ; \(\overline{f}\) is an extension of \(f\) since, if \(x \in A\) , all the components \(x_i\) of \(x\) on the lines \(D_i\) also belong to \(A\) , by reason of the choice of the \(a_i\) ; moreover, \(\overline{f}\) is continuous on \(R^n\) , because it is continuous on each of the lines \(D_i\) , and \(x_i\) is a linear function of \(x\) (and therefore continuous).

COROLLARY 1. Let V be a neighbourhood of o in \(R^{n}\) and let f be a continuous mapping of V into a topological group G such that

\[
f (\boldsymbol {x} + \boldsymbol {y}) = f (\boldsymbol {x}) f (\boldsymbol {y})
\]

for each pair of points \(x, y\) such that \(x \in V\) , \(y \in V\) and \(x + y \in V\) . Then there exists a unique continuous homomorphism of \(\mathbf{R}^n\) into \(G\) which agrees with \(f\) at all points of a neighbourhood \(W\) of \(o\) .

Take W to be an open box with centre o, contained in V, and apply Proposition 2 to W.

We shall see later that this property of \(R^{n}\) extends to a larger class of topological groups, the ``simply connected'' groups.

COROLLARY 2. Let \(f\) be a local isomorphism of \(\mathbf{R}^n\) with a topological group \(\mathbf{G}\) . Then there exists a unique strict morphism of \(\mathbf{R}^n\) onto an open subgroup of \(\mathbf{G}\) which agrees with \(f\) at all the points of some neighbourhood of \(\mathbf{o}\) .

Let \(\overline{f}\) be the continuous homomorphism of \(\mathbf{R}^n\) into \(\mathbf{G}\) which agrees with \(f\) at all points of a neighbourhood of o; \(\overline{f}(\mathbf{R}^n)\) contains, by hypothesis, a neighbourhood of the identity element of \(\mathbf{G}\) and therefore (Chapter III, § 2, no. 1, Corollary to Proposition 4) is an open subgroup of \(\mathbf{G}\) ; moreover \(\overline{f}\) is a strict morphism of \(\mathbf{R}^n\) onto \(\overline{f}(\mathbf{R}^n)\) , by Chapter III, § 2, no. 8, Proposition 24.

THEOREM 1. Every connected group G which is locally isomorphic to \(\mathbf{R}^n\) is isomorphic to a group \(\mathbf{R}^p \times \mathbf{T}^{n-p}\) ( \(0 \leqslant p \leqslant n\) ).

A suitably chosen local isomorphism \(f\) of \(\mathbf{R}^n\) with \(G\) extends to a strict morphism of \(\mathbf{R}^n\) onto an open subgroup of \(G\) (Proposition 2, Corollary 2) and therefore onto \(G\) itself, since \(G\) is connected. It follows that \(G\) is isomorphic to a quotient group \(\mathbf{R}^n / H\) of \(\mathbf{R}^n\) : \(H\) is discrete, otherwise there would exist points \(x \neq 0\) of \(H\) arbitrarily close to \(o\) and such that \(f(x) = f(0)\) , contrary to the hypothesis that \(f\) is a local isomorphism. The theorem is therefore a consequence of Theorem 1 of § 1, no. 1.

